\documentclass[11pt]{article}

\usepackage[margin=1in]{geometry}
\usepackage{amsmath,amssymb}
\usepackage{booktabs}
\usepackage{hyperref}
\hypersetup{hidelinks}

\title{Scaling Book Exercises: Section 5}
\author{Lucas Yan}
\date{}

\begin{document}
\maketitle

\noindent
\textbf{Chapter:} How to Parallelize a Transformer for Training\\
\textbf{Source scan:} \texttt{../handwritten/section\_05\_handwritten.pdf}\\
\textbf{Reference:} \url{https://jax-ml.github.io/scaling-book/training/}

\paragraph{Model and hardware.}
The worked problems use LLaMA-2 13B with
\[
L=40,\quad D=5120,\quad F=13824,\quad N=K=40,\quad
H=128,\quad V=32000.
\]
The input and output embeddings are separate. TPU v5p has 96 GB of HBM,
peak bfloat16 compute
\[
C=4.59\times10^{14}\ \text{FLOPs/s},
\]
and an ICI operational intensity of approximately 2550.

\section*{Question 1: Parameter count}

The gated MLP contains three large matrices per layer:
\[
P_{\mathrm{MLP}}=3LDF
=3(40)(5120)(13824)
=8{,}493{,}465{,}600.
\]
Because this model has the same number of query and KV heads, attention has
four $D\times NH$ projections per layer:
\[
P_{\mathrm{attention}}=4LDNH
=4(40)(5120)(40)(128)
=4{,}194{,}304{,}000.
\]
The two vocabulary matrices contribute
\[
P_{\mathrm{vocab}}=2VD
=2(32000)(5120)
=327{,}680{,}000.
\]
Therefore,
\begin{align*}
P
&=P_{\mathrm{MLP}}+P_{\mathrm{attention}}+P_{\mathrm{vocab}}\\
&=13{,}015{,}449{,}600\\
&\approx\boxed{13.0\text{B parameters}}.
\end{align*}

\section*{Question 2: Training memory at a 16M-token batch}

With bfloat16 parameters and two float32 Adam accumulators, each parameter
requires
\[
2+4+4=10\ \text{bytes}.
\]
Thus the parameters and optimizer state occupy
\[
M_{\mathrm{model}}=10P
=130.15\times10^9\ \text{bytes}
\approx130\ \text{GB}.
\]

The three checkpointed activations per layer have shapes $[B,F]$, $[B,F]$,
and $[B,D]$. In bfloat16, their total size is
\begin{align*}
M_{\mathrm{activations}}
&=2LB(D+2F)\\
&=2(40)(16\times10^6)(5120+2(13824))\\
&=41.943\times10^{12}\ \text{bytes}\\
&\approx41.94\ \text{TB}.
\end{align*}
The total training state is therefore approximately
\[
\boxed{42.07\ \text{TB}},
\]
with activation checkpoints dominating the memory requirement.

\section*{Question 3: A 4096-chip TPU v5p slice}

The setting is a $16\times16\times16$ slice, so
\[
N_{\mathrm{chips}}=4096,
\]
with a 3M-token batch and 32k sequence length. This is roughly 96 sequences
per step.

\subsection*{(a) Pure data parallelism}

Pure data parallelism replicates the full parameters and optimizer state on
every chip. The replicated state alone is about 130 GB, exceeding the 96 GB
HBM capacity of one TPU v5p chip. Therefore,
\[
\boxed{\text{pure data parallelism does not fit in memory}.}
\]

\subsection*{(b) Pure FSDP}

Replacing the 16M-token batch in Question 2 with 3M tokens gives
\begin{align*}
M_{\mathrm{activations}}
&=2(40)(3\times10^6)(5120+2(13824))\\
&=7.864\ \text{TB}.
\end{align*}
Including parameters and optimizer state gives approximately 7.99 TB total,
or
\[
\frac{7.99\ \text{TB}}{4096}
\approx\boxed{1.95\ \text{GB per chip}}.
\]
Thus FSDP easily fits in aggregate memory.

The issue is communication. Using all three mesh axes for FSDP, the minimum
compute-bound local batch is
\[
\frac{2550}{3}=850\ \text{tokens per chip}.
\]
The actual local batch is
\[
\frac{3\times10^6}{4096}\approx732\ \text{tokens per chip}.
\]
Since $732<850$, pure FSDP is feasible but communication-bound. It does not
remain in the strong-scaling regime.

\subsection*{(c) Mixed FSDP and tensor parallelism}

For mixed FSDP and tensor parallelism, the compute-bound requirement is
\[
\frac{B}{N}>
\frac{2550^2}{2F}
=\frac{2550^2}{2(13824)}
\approx235\ \text{tokens per chip}.
\]
The actual 732 tokens per chip clears this threshold.

Let $X$ be the FSDP degree and $Y$ the tensor-parallel degree, with $XY=N$.
Using two mesh axes for FSDP and one for tensor parallelism,
\begin{align*}
X_{\mathrm{opt}}
&=\sqrt{\frac{B}{F}\frac{M_X}{M_Y}N}\\
&=\sqrt{\frac{3\times10^6}{13824}(2)(4096)}\\
&\approx1333.
\end{align*}
A practical power-of-two choice is
\[
\boxed{X=1024,\qquad Y=4}.
\]

Using the standard training estimate of six FLOPs per parameter per token,
the ideal step time at 40\% MFU is
\begin{align*}
T_{\mathrm{step}}
&=\frac{6PB}{N(0.4)C}\\
&=\frac{6(13.015\times10^9)(3\times10^6)}
{4096(0.4)(4.59\times10^{14})}\\
&\approx\boxed{0.31\ \text{s per step}}.
\end{align*}

\end{document}
